\chapter{Threat Model}\label{threat:chapter}

\section{Method and framing}\label{threat:method}

The preceding chapters presented the five mechanisms of CAS Core 8 constructively: semantic references and typed links, the research boundary object, the canonical evidence encoding, the compatibility worker protocol with its concurrency and failure disciplines, and the byte-backed registries feeding the runtime decision model. This chapter inverts the presentation. Instead of asking what each mechanism does, it asks what each mechanism is \emph{for}: which adversarial or accidental behavior, if unblocked, would corrupt which artifact, and what happens in the kernel when the behavior is attempted.

The method is deliberately conventional. We follow the threat-modeling discipline of enumerating assets, trust boundaries, and per-threat mitigations \cite{shostack}, and the broader security-engineering stance that a system's real specification is the set of failures it survives rather than the set of features it advertises \cite{anderson}. Two adaptations are worth stating up front, because they shape every entry below.

First, the assets of this kernel are not user data or money; they are \emph{evidence artifacts}: semantic references and the typed links between them, boundary objects, canonical encodings and their digests, compatibility receipts, registry snapshots, and the decision records that consume all of the above. An attack succeeds when a consumer of one of these artifacts is entitled, by the kernel's own rules, to draw a conclusion the artifact does not support. This is a weaker success criterion than code execution or data exfiltration, and it is the correct one for a research program whose historically observed failures were precisely conclusions drawn from artifacts that did not support them: positional relabels read as paraphrases, and structural receipts read as licenses.

Second, because the kernel's design stance is fail-closed rather than detect-and-respond, the canonical mitigation shape here is \emph{unrepresentability} or \emph{rejection at construction}, not monitoring. Where a classical threat model would say ``log and alert,'' this kernel says ``the state cannot be constructed'' or ``the loader raises.'' The residual-risk paragraphs below are therefore where the honest content lives: they state what remains possible when the representable-state restriction is not the whole story.

\begin{definition}[Per-threat treatment format]\label{threat:def:format}
Each threat in this chapter is treated in five fixed parts: (i) a \emph{scenario narrative} describing a concrete attempted corruption; (ii) the \emph{preconditions} an attacker or accident requires; (iii) the \emph{artifact corrupted} if the threat succeeds, named from the asset list of Section~\ref{threat:assets}; (iv) the \emph{blocking mechanism}, citing the kernel rule or protocol step that stops it; and (v) the \emph{residual risk} that the blocking mechanism does not remove.
\end{definition}

\begin{remark}\label{threat:rem:attacker}
The attacker model deserves explicitness. The kernel defends against: (a) code running in the same Python process as the kernel, including well-meaning tooling and AI-authored agents, that supplies crafted inputs, subclasses public types, or constructs objects through non-loader paths; (b) suppliers of files the kernel reads --- source programs, registry ledgers, schema files, compiler bytes; (c) concurrent processes contending for the compatibility worker; and (d) drift over time --- interpreter upgrades, dependency changes, platform newline conventions --- which behaves adversarially even without an adversary. The kernel does \emph{not} defend against an attacker with operating-system privileges, control of native libraries, or control of the trust roots that a future authorization design would introduce; Section~\ref{threat:scope} treats these exclusions honestly. Within this model, note the historical fact recorded in the program's own audit trail: several of the incidents motivating this chapter were caused by AI-authored tooling acting without malice. The model therefore does not distinguish malice from enthusiasm; both produce the same corrupted artifacts.
\end{remark}

\section{Assets and trust boundaries}\label{threat:assets}

The artifacts at risk, in dependency order:

\begin{enumerate}
\item \emph{Semantic references and typed links.} The structural tuples of the form (namespace, \texttt{local\_id}, version, \texttt{source\_id}, \texttt{definition\_sha256}) and the \textsc{translation} / \textsc{alignment} / \textsc{hypothesis} links between them. Corruption here poisons every downstream statement about what a piece of CAS content \emph{means}.
\item \emph{Boundary objects.} The (admission, product, disclosure) triples. Corruption here converts research material into apparently licensed product.
\item \emph{Canonical encodings and digests.} The deterministic byte forms over which all hash equality is computed. Corruption here silently invalidates every comparison in the kernel at once.
\item \emph{Compatibility receipts.} The records binding source identity, compiler identity, worker runtime, engine profile, and result digest. Corruption here manufactures false compatibility evidence.
\item \emph{Registry snapshots and selected records.} The byte-backed claim and theorem registries. Corruption here rewrites the program's claim history.
\item \emph{Decision records.} The outputs of the runtime decision model. Corruption here is the terminal failure: an action proceeds under a justification the kernel never issued.
\end{enumerate}

The trust boundaries are: the public constructor surface of every kernel type (inside which exact-type and validity checks run); the loader path of every registry (the only admitted route from bytes to records); the worker process boundary of the compatibility profile; and the filesystem, which the kernel reads but never trusts without hashing.

With the format, attacker model, and assets fixed, we treat eleven threats.

\section{T1: Namespace conflation}\label{threat:t1}

\begin{threat}[Positional correspondence laundered to identity]\label{threat:thr:conflation}
An agent treats a positional or historical correspondence between elements of different ontological families as an identity, and conclusions established for one family propagate to the other without any record that a crossing occurred.
\end{threat}

\paragraph{Scenario.} A tool ingests documents from two eras of the program. One names \texttt{F5} as a TKS Foundation position; another names \texttt{F5} as a CAS Drive. The strings are equal; the concepts are related by position and history but are not interchangeable, and prior audit found exactly this conflation voiding a completed work product. The tool builds a lookup keyed by the bare string, merges the two families into one table, and every downstream consumer now silently inherits statements about Foundations as if they were statements about Drives. No step in the pipeline is individually wrong; the corruption is in the key.

\paragraph{Preconditions.} Untyped string identifiers anywhere on the path from source text to conclusion; a correspondence that is real (so the merge looks justified) but weaker than identity; and no forced representation of link kind. The threat requires no malice --- it is the default behavior of any dictionary keyed on rendered names.

\paragraph{Artifact corrupted.} Semantic references and links (asset 1), and transitively everything built on them: an analysis, a claim record citing the analysis, a decision consuming the claim.

\paragraph{Blocking mechanism.} The semantic-reference definition assigns TKS Foundations, TKS Worlds, CAS Drives, CAS Domains, CAS Registers, CAS States, and Noetic Topos objects to disjoint namespaces, and identity is structural equality of the \emph{full} tuple including namespace and \texttt{definition\_sha256}. Two references with equal local identifiers in different namespaces are unequal by construction. The no-silent-identity rule then closes the remaining route: a cross-namespace correspondence is representable only as a typed link of kind \textsc{translation}, \textsc{alignment}, or \textsc{hypothesis}, each carrying its own grounding-artifact digest, mandatory note, and Ring-4 boundary, and there is no coercion path from any link kind to identity. The strongest statement an agent can construct is ``these correspond, per this source, at this review status'' --- which is exactly the statement the historical audit said should have been made.

\paragraph{Residual risk.} The kernel governs its own representations, not free-text prose. An agent that never constructs semantic references --- that reasons entirely in untyped strings outside the kernel --- can still conflate; the kernel's protection is only as broad as its adoption. Second, a \textsc{translation} link honestly constructed can still be over-read by a consumer who ignores its kind field; the kernel makes the kind unremovable, not unignorable. Third, the disjointness of namespaces is itself a modeling decision: if a future ontology genuinely merges two families, the kernel will require an explicit, reviewed representation of that merge rather than permitting it, which is the intended bias but is still a cost.

\section{T2: Identifier collision via delimiter injection}\label{threat:t2}

\begin{threat}[Crafted components alias two references to one rendered identifier]\label{threat:thr:collision}
Components containing delimiter characters are chosen so that two distinct structural tuples render to the same identifier string, aliasing their records everywhere the rendered form is used as a key.
\end{threat}

\paragraph{Scenario.} Rendered identifiers join tuple components with a delimiter, say a colon. An input source defines a local identifier that itself contains a colon, chosen so that (\texttt{ns}, \texttt{a:b}, \ldots) renders identically to (\texttt{ns:a}, \texttt{b}, \ldots). A log index, a cache, or a registry keyed on rendered identifiers now merges the two references' histories. The attack is the classic in-band signaling failure: control characters travel in the data channel. Its accidental form is equally realistic --- legacy CAS notation and file-derived source identifiers contain punctuation freely.

\paragraph{Preconditions.} A rendering that concatenates without escaping or length-prefixing; at least one component alphabet that includes the delimiter; and a consumer that keys on the rendered string rather than the structural tuple.

\paragraph{Artifact corrupted.} Semantic references as they appear in any keyed context (asset 1), plus any evidence record whose canonical encoding embeds the rendered form.

\paragraph{Blocking mechanism.} The injective-rendering invariant: rendering escapes or length-prefixes every component before joining, so the rendered form is uniquely parseable and the map from structural tuples to rendered strings is injective. Two distinct tuples differ at the first divergent component under any parse of their renderings. The property is enforced by construction, and the Gate-4 suite regression-tests it with delimiter-bearing component fixtures --- components that contain the delimiter, that are empty, and that are prefixes of one another.

\paragraph{Residual risk.} Injectivity is a property of the kernel's renderer, not of every renderer. External systems that re-render references through their own concatenation --- a spreadsheet export, a log formatter --- can reintroduce the collision downstream of the kernel; the mitigation available to them is to carry the structural tuple or the kernel's rendered form verbatim, never to re-derive it. There is also a subtler residue: injectivity guarantees distinct tuples render distinctly, but near-collisions (identifiers differing by visually confusable characters) remain possible in principle. The identifier-profile rule confines review-visible machine identifiers to declared normalized/ASCII profiles precisely to shrink this residue, but human misreading of two honestly distinct identifiers is not a machine-checkable failure.

\section{T3: Equality spoofing via subclass override}\label{threat:t3}

\begin{threat}[A subclass overrides equality to impersonate an identity-bearing value]\label{threat:thr:spoof}
A class derived from a public kernel type overrides its equality or hashing behavior so that an impostor instance compares equal to an honest reference or boundary object, passing every check that compares values.
\end{threat}

\paragraph{Scenario.} Python's object model makes this threat nearly free to attempt. A tool defines \texttt{class BoundaryEx(ResearchBoundary)} with \texttt{\_\_eq\_\_} returning \texttt{True} against any boundary, or a reference subclass whose equality ignores \texttt{definition\_sha256}. The impostor is then handed to a kernel operation that checks ``is this the Ring-4 boundary?'' by comparison. Every comparison-based guard in the process now answers yes to an object whose actual field values say otherwise. The scenario needs no privilege beyond the ability to define a class in the same interpreter --- which every plugin, notebook, and AI-authored helper has.

\paragraph{Preconditions.} A kernel that accepts subclass instances at its public interfaces, or that implements guards via symmetric equality where the impostor's method can win dispatch.

\paragraph{Artifact corrupted.} Boundary objects and semantic references at the point of use (assets 1 and 2): not the stored artifacts themselves, but every decision predicated on comparing them.

\paragraph{Blocking mechanism.} The exact-type identity rule: equality-overriding subclasses are not admitted as identity-bearing public values. Kernel checks compare exact runtime types first --- \texttt{type(x) is} the expected class, not an \texttt{isinstance} test --- and then primitive fields, so an impostor is rejected before its equality method can participate. Public constructors likewise reject subclass instances of their argument types. The check is deliberately hostile to a common object-oriented idiom; the kernel's position is that identity-bearing values are closed data, not extension points.

\paragraph{Residual risk.} Exact-type checking protects the kernel's own comparisons. Code outside the kernel that compares kernel values with bare \texttt{==} remains spoofable within that code's own logic; the kernel cannot police foreign comparisons. Additionally, exact-type checks are a Python-level defense: an attacker who can mutate the class object itself, patch the kernel's modules, or interfere below the interpreter is outside this mechanism's reach and inside the excluded scope of Section~\ref{threat:scope}. Finally, the rule imposes a real engineering cost --- no convenient subclassing of kernel types even for benign wrappers --- and the pressure to relax it is itself a standing risk that only review discipline addresses.

\section{T4: Digest instability and type puns}\label{threat:t4}

\begin{threat}[Nondeterministic or ambiguous encoding breaks hash equality in both directions]\label{threat:thr:digest}
Equal evidence hashes differently across runs or platforms (false difference), or different evidence hashes equally (false identity), because the encoding admits iteration-order nondeterminism, float-formatting variance, or cross-type collisions among $\mathrm{Decimal}$, string, and integer values.
\end{threat}

\paragraph{Scenario.} Three concrete instances. (i) A record is encoded from a hash-ordered container; two runs serialize its keys in different orders and the same evidence carries two digests, so replay comparison fails against an honest re-execution --- or worse, a comparison framework ``helpfully'' relaxes to set-comparison and stops detecting real differences. (ii) A measured quantity passes through a binary float; \texttt{0.1 + 0.2} formats differently across serializer versions, and a receipt's result digest drifts with no semantic change. NaN payloads and infinities serialize non-portably or compare unequal to themselves, poisoning any structure containing them. (iii) A value recorded as $\mathrm{Decimal}(\texttt{"1.0"})$ in one producer and as the string \texttt{"1.0"} or the integer $1$ in another encodes to the same bytes; two records asserting different typed facts become digest-identical, and a type pun becomes an evidence forgery primitive: construct a benign record whose encoding collides with a meaningful one.

\paragraph{Preconditions.} Any encoder on the evidence path that tolerates unordered containers, binary floating point, or untagged scalar coercion; or more than one serialization route for the same public type, so that two honest producers disagree byte-wise.

\paragraph{Artifact corrupted.} Canonical encodings and digests (asset 3) --- and because every registry record, compatibility receipt, and replay comparison reduces to hash equality of canonical encodings, an instability here invalidates assets 4 through 6 wholesale. This is the highest-leverage threat in the model: it does not corrupt one artifact but the comparison relation itself.

\paragraph{Blocking mechanism.} The canonical-encoding rule set, specified closed: deterministic UTF-8 JSON with sorted string keys, in the spirit of RFC 8785 canonicalization \cite{jcs}; outright rejection --- not normalization --- of sets and other unordered containers, binary floats, NaN, and infinity; tagged exact decimals so that $\mathrm{Decimal}(\texttt{"1.0"})$, the string \texttt{"1.0"}, and the integer $1$ have pairwise distinct encodings; exactly one declared serialization route per public type; and bounded, cycle-aware canonicalization that fails closed on cycles and on graphs exceeding declared depth or node limits. Rejection rather than coercion is the load-bearing choice: a float is not rounded into acceptability, it is an error at the boundary where the producer can still fix it.

\paragraph{Residual risk.} Determinism is proven only over the value domain the encoder admits; a future extension of the domain (a new tagged type) reopens the analysis and must arrive with its own collision argument. The scheme also inherits the collision resistance of SHA-256: digest equality is treated as byte equality, which is a standard but stated cryptographic assumption \cite{anderson}. And boundedness has a denial aspect --- an honest but enormous evidence graph fails closed; the declared limits are a policy choice that trades availability for the guarantee that canonicalization terminates. Finally, the one-route-per-type rule is enforced for kernel types only; foreign objects reach the encoder only by prior conversion, and a sloppy conversion layer can reintroduce ambiguity one step before the boundary.

\section{T5: Newline rewriting}\label{threat:t5}

\begin{threat}[Platform newline translation masks a real difference or fakes a spurious one]\label{threat:thr:newline}
CRLF/LF rewriting between checkout, editing, and transport changes raw bytes without semantic change (spurious difference), or a normalization step erases a genuine byte-level difference before it is recorded (masked difference).
\end{threat}

\paragraph{Scenario.} The program runs on Windows hosts with Git newline translation in the loop; the failure is not hypothetical but part of the kernel's own history, where a line-ending repair episode had to be frozen and receipted. In one direction: a compiler source file checked out with CRLF hashes differently from the LF bytes a receipt recorded, and an honest replay fails identity verification --- pressure then builds to ``just normalize,'' which is the second direction: a normalizing comparison silently equates two files that differ by a lone CR embedded mid-line, a difference that may be semantically meaningful to a parser and is certainly meaningful to byte identity. A single-digest design must choose one failure or the other.

\paragraph{Preconditions.} Cross-platform file transport; any single-identity design that either always normalizes or never normalizes; and tooling (editors, VCS filters) that rewrites silently.

\paragraph{Artifact corrupted.} Compatibility receipts and any byte-identity field (assets 3 and 4): either replayability is destroyed by spurious mismatch, or tamper evidence is destroyed by over-normalization.

\paragraph{Blocking mechanism.} The byte-identity-versus-comparison-identity rule: raw bytes and the declared CRLF-to-LF comparison identity are recorded as \emph{separate fields with distinct roles}, and lone-CR, CRLF, and lone-LF counts are recorded rather than conflated. A receipt therefore carries both the raw digest (tamper evidence, exact provenance) and the LF-normalized digest (declared comparison identity), and neither substitutes for the other. A platform rewrite now surfaces as: raw digests differ, comparison digests match, newline counts explain the delta --- a fully diagnosed, honest state rather than either failure mode. The compatibility protocol records raw and LF compiler identities side by side for exactly this reason.

\paragraph{Residual risk.} The comparison identity is CRLF-to-LF by declaration; other whitespace equivalences (trailing spaces, BOM presence, tab/space) are \emph{not} normalized, and a rewrite of those still produces a raw-and-comparison double mismatch that must be adjudicated by a human. This is intended --- the kernel refuses to guess which equivalences are semantic --- but it means newline discipline does not generalize to a full ``semantic identity'' for text, and no such identity is claimed. The counts also diagnose only newline rewrites; they do not identify \emph{which} tool rewrote, which remains an environment-forensics question outside the receipt.

\section{T6: Compiler-identity drift}\label{threat:t6}

\begin{threat}[A receipt is produced against different compiler bytes than claimed, or under inherited interpreter state]\label{threat:thr:compdrift}
A compatibility execution runs a compiler that is not byte-identical to the one the receipt names --- because the file changed between hash and use, because import resolution found a different module, or because the parent interpreter's loaded state leaked into the run --- and the receipt binds a result to an identity that did not produce it.
\end{threat}

\paragraph{Scenario.} Three variants. (i) Time-of-check to time-of-use: the compiler file is hashed, then edited (by a build step, a sync client, an agent) before execution; the receipt's compiler digest is honest about the past and wrong about the run. (ii) Import shadowing: the execution imports the compiler by module name and the interpreter resolves a different file --- a stale \texttt{.pyc}, a path-shadowing local module, a site customization --- so the hashed bytes and the executed bytes were never the same object. (iii) Inherited state: the run executes in the parent process, where dozens of already-imported modules, monkeypatches, and environment variables shape behavior; the same compiler bytes yield different results across parents, and the receipt has no field that could express why.

\paragraph{Preconditions.} Any gap between hashing and execution that the filesystem can move through; name-based rather than byte-based execution; or in-process execution inheriting \texttt{sys.modules}, site configuration, or environment.

\paragraph{Artifact corrupted.} Compatibility receipts (asset 4): specifically the binding between compiler identity and result digest, which is the receipt's entire value.

\paragraph{Blocking mechanism.} Deterministic-closure protocol steps taken together. Step (1) verifies the requested source hash and reads and hashes the compiler bytes \emph{before} execution; step (4) then executes those verified bytes --- not a module name --- in a fresh \texttt{python -I -S -B} worker with a sanitized environment, so isolated mode, no site imports, no bytecode caching, and no inherited \texttt{sys.modules} remove the shadowing and inherited-state variants; step (5) samples and binds the worker runtime (interpreter fields, \texttt{sys.platform} or POSIX \texttt{os.uname}, profile-bound Windows architecture inputs, without WMI or subprocess probing) before compiler code executes, so the receipt records the actual runtime rather than an assumption; and step (7) records raw and LF compiler identities, interpreter and worker bytes, import search paths, and the loaded-module inventory, making any residual leakage inspectable. The replay requirement of step (8) --- treated separately as T9 --- closes the temporal gap: a structural decision requires a fresh execution against the same source and compiler identity, so a one-time lucky binding cannot be reused indefinitely.

\paragraph{Residual risk.} The hash-then-execute window inside a single run is minimized, not eliminated: the kernel reads bytes and passes those bytes to the worker, but the interpreter binary itself, and every native library it loads, are identified rather than content-addressed --- the stated non-sandbox boundary. A host that swaps the Python executable or a shared library between receipt and replay defeats the binding at a layer the protocol only records, does not verify. This is exactly open obligation one (hermetic bundle) and is deliberately surfaced in every receipt as \texttt{false} capability fields rather than hidden.

\section{T7: Receipt laundering}\label{threat:t7}

\begin{threat}[Structural receipts or self-hashes are presented as authority]\label{threat:thr:launder}
An artifact that proves at most structural consistency --- a green test run, a self-computed digest, a generated document with the word ``licensed'' in it --- is presented to a human or downstream system as if it established reviewer identity, legal authority, or empirical validity.
\end{threat}

\paragraph{Scenario.} This threat has the strongest historical grounding in the model, and stating it plainly is part of the mitigation. In the program's own past, AI-authored tooling generated ``licensing receipts'': documents asserting licensed status, carrying self-computed hashes of their own content, formatted with the visual apparatus of authority. Every property of those artifacts was structurally sound --- the hashes verified, the formats parsed --- and none of it constituted authority, because no authenticated reviewer, admission process, or legal trust root had touched them. The self-hash answered ``has this document changed since it was written?'' and was read as answering ``who stands behind this document?''. The laundering step is precisely that category slide, and it required no forgery: the artifacts were genuine artifacts of the wrong kind. The general form recurs constantly in software practice: CI badges read as safety cases, checksums read as signatures, passing suites read as review.

\paragraph{Preconditions.} A representable ``authorized'' state that structural evidence can reach; consumers who do not distinguish tamper evidence from provenance; and production of official-looking artifacts being cheap --- which, with generative tooling, it now permanently is.

\paragraph{Artifact corrupted.} Decision records and, socially, the program's claim discipline itself (asset 6): the terminal corruption is an action taken, or a status asserted, on the strength of an artifact that cannot bear it.

\paragraph{Blocking mechanism.} Three layers, converging on unrepresentability. The recheck invariant: every public operation derives and rechecks the research boundary from its inputs, never trusting a shared mutable object or a recomputable self-hash, and the kernel states explicitly that a self-digest is tamper evidence, not signature, reviewer identity, or authority grant. The metadata-is-not-authority invariant: the text \textsc{licensed} inside a registry record is necessary metadata, never sufficient authority. And terminally, the no-authorization-path invariant: in v0.1 the field \texttt{authorization\_granted} and the compatibility spelling \texttt{allowed} are constant \emph{false}, with no code path that sets either true; claim-bearing intent is denied even when a locally constructed ledger says \textsc{licensed}. There is, in this kernel, no authorized state to launder \emph{into}. A laundered receipt can deceive a human reading prose; it cannot cause the kernel to emit a decision record that says yes.

\paragraph{Residual risk.} The kernel controls its own outputs, not human interpretation of them. A person can still over-read a Gate-4 receipt in a slide deck; the mitigation there is the claim-discipline language this monograph itself practices --- receipts are historical software evidence under exact fixtures, Ring 4, \textsc{not licensed} --- and that mitigation is social, not mechanical. The constant-\emph{false} invariant also has a planned expiry: a future authorization design will make ``yes'' representable under the constraints stated in the runtime decision model (authenticated authority and trust root, reviewer and admission receipts, exact registry manifest, bound operation identities). The laundering threat then returns in full force against that design, which is why its constraint list is stated now and why independent review gates any status change. Until then, the residual risk of T7 inside the kernel is, unusually, zero by construction --- and outside the kernel, unchanged.

\section{T8: Registry substitution}\label{threat:t8}

\begin{threat}[Parsed content is swapped under a retained raw identity]\label{threat:thr:substitution}
A registry object presents the raw-byte identity of an honest ledger while its parsed records have been replaced, edited, or reordered --- via direct construction, dataclass copying, or mutation after load --- so consumers verify the right bytes and read the wrong content.
\end{threat}

\paragraph{Scenario.} The registry API exposes, as it must, both the raw identity (digests of ledger and schema bytes) and the parsed records. An in-process tool constructs a registry object directly --- calling the record dataclass constructor, or copying a loaded object and editing a field --- pairing the honest raw digests with altered content: a claim record's ring changed, a version bumped, a record deleted. Downstream code that ``verifies the registry'' by checking the raw digest passes, then consumes the substituted parse. The threat also has an accidental form: a caching layer that rehydrates records from its own serialization, drifting from the ledger bytes it claims to represent.

\paragraph{Preconditions.} A construction path from parsed values to a registry object that does not pass through the bytes; or post-load mutability of records; or record digests computed over parsed values alone, unbound to the ledger bytes.

\paragraph{Artifact corrupted.} Registry snapshots and selected records (asset 5), and through them the decision model's claim inputs.

\paragraph{Blocking mechanism.} The byte-backed-loading rule. The loader is the only admitted route from bytes to records: parsed headers, schemas, and records are reconstructed from and rechecked against the retained raw ledger and schema bytes, so a direct constructor or dataclass copy cannot preserve an old raw identity while substituting parsed content --- the recomputation from bytes disagrees and the object is rejected. Loaded values are recursively immutable, closing the post-load mutation route. And each selected record carries a deterministic digest bound to the registry bytes, not merely to its own field values, so a record cannot be excised from one ledger context and presented under another. The loader's rejection catalog --- duplicate JSON keys, unknown root or record fields, wrong enums or shapes, invalid identifiers, ambiguous versions, missing required metadata --- additionally removes the parser-differential route, where two consumers parse the same bytes into different contents: bytes that are ambiguous under any of those headings never load at all.

\paragraph{Residual risk.} Byte-backing authenticates content against bytes, not bytes against the world: an attacker who can rewrite the ledger file itself presents self-consistent substituted bytes, and the loader will honestly load them. What blocks that step is not the loader but the digest pinning upstream --- receipts and reviews that name expected ledger digests --- and ultimately the trust-root question deferred to future authorization design. Within the process, Python's mutability escape hatches (\texttt{object.\_\_setattr\_\_}, direct \texttt{\_\_dict\_\_} manipulation) can in principle pierce recursive immutability; this sits at the same interpreter-integrity layer as T3's residue and is excluded scope. The practical residue is thus: the loader guarantees you read what the bytes say; something else must guarantee they are the right bytes.

\section{T9: Replay evasion}\label{threat:t9}

\begin{threat}[A structural decision is issued without fresh replay]\label{threat:thr:replay}
A stored compatibility receipt is accepted as sufficient for a structural decision without re-executing against the same source and compiler identity, so decisions ride on evidence whose validity conditions may have silently lapsed.
\end{threat}

\paragraph{Scenario.} A receipt from last month verifies perfectly: digests intact, boundary consistent, all fields well-formed. A consumer, reasonably minded to save ninety seconds, treats verification of the stored receipt as equivalent to current compatibility and issues a structural pass. Meanwhile the interpreter was patched, a dependency wheel was rebuilt, or the compiler file was touched and reverted --- none of which the stored receipt can reflect, because it is a statement about a past execution. The gap between ``this receipt was true'' and ``this is true now'' is exactly the gap replay exists to close; evading replay converts historical evidence into present authority, a temporal cousin of T7.

\paragraph{Preconditions.} A decision path that accepts receipt verification in place of re-execution; caching layers with unbounded or unprincipled lifetimes; or cost pressure institutionalizing ``verify the artifact, skip the run.''

\paragraph{Artifact corrupted.} Decision records (asset 6): a structural pass whose stated basis --- current compatibility under the named identities --- did not hold at decision time.

\paragraph{Blocking mechanism.} Step (8) of the deterministic-closure protocol is a hard gate: fresh replay against the same source and compiler identity is required \emph{before any structural decision can pass}. The decision model has no input slot through which a stored receipt alone can satisfy the compatibility predicate --- what it consumes is the agreement between a stored identity and a fresh execution's identity and result digest. Because the fresh run re-verifies source and compiler hashes (step 1) and re-samples the worker runtime (step 5), the lapsed-validity scenarios surface as identity mismatches or result-digest disagreements rather than as silent staleness. Replay is thus not a redundancy check but the mechanism by which time re-enters the evidence.

\paragraph{Residual risk.} Replay establishes that the execution reproduces \emph{now}, under the same recorded identities --- it does not establish determinism across environments the protocol does not bind (see T6 residue and T11), and a quantity that is stable on one host may still diverge elsewhere. Replay also doubles execution cost by design; the predictable institutional pressure to reintroduce caching is a governance risk the kernel answers only by making the gate structural rather than advisory. Finally, replay freshness is bounded by the granularity of a run: an environment change \emph{during} the replay window is theoretically unobserved, though the worker's short lifetime and pre-execution runtime sampling make the window small.

\section{T10: Lock-based denial}\label{threat:t10}

\begin{threat}[The advisory lock is held to deny compatibility executions]\label{threat:thr:lock}
A process acquires the profile-bound advisory lock and holds it --- maliciously, or by crashing at the wrong moment --- so that cooperating workers cannot execute and compatibility evidence cannot be produced.
\end{threat}

\paragraph{Scenario.} The compatibility profile serializes cooperating worker executions through a profile-bound advisory lock. A stuck process --- a debugger session left at a breakpoint inside the critical region, a crashed worker whose lock artifact lingers, or a deliberately hostile loop that reacquires immediately after release --- starves other workers. No evidence is corrupted; evidence is \emph{prevented}, and a program that must produce receipts on a cadence experiences this as an outage.

\paragraph{Preconditions.} Shared-host contention for the same profile lock; a holder that does not release promptly; and, for the starvation variant, scheduling that favors the reacquiring process.

\paragraph{Artifact corrupted.} None --- this is the model's one pure availability threat, and its inclusion is precisely to record that the integrity mechanisms do not convert it into anything worse.

\paragraph{Blocking mechanism.} Bounded impact rather than prevention. Acquisition is bounded at 90 seconds with a 10-millisecond poll interval, so a denied waiter fails \emph{closed and diagnosably} --- a timeout error, not a hang and not a lockless fallback execution. A 10-millisecond post-release scheduler yield reduces (without eliminating) immediate-reacquisition starvation. Two honesty clauses bound the claim itself: a failed worker is never retried, so denial cannot be silently converted into selection bias over receipts (the T11-adjacent hazard of retry-until-quiet); and the protocol explicitly does \emph{not} claim lock fairness, so no consumer can build on an ordering guarantee the mechanism does not provide. The worst case is therefore stated and small: no evidence for the duration of the denial, never wrong evidence because of it.

\paragraph{Residual risk.} The full availability residue: a determined local adversary can deny compatibility execution indefinitely, and the kernel's only answer is that the failure is loud. Advisory locks bind only cooperating processes --- a non-cooperating process does not corrupt the lock discipline's guarantees among cooperators, but two runs that ignore the profile lock entirely are outside its serialization, relying on the worker protocol's per-run isolation instead. Timeout tuning is a policy residue: 90 seconds is a declared constant, not a derived optimum, and hosts with legitimately long compiler runs under contention will see spurious timeouts before they see starvation.

\section{T11: Dependency-profile divergence between workers}\label{threat:t11}

\begin{threat}[Two successful workers disagree on their dependency profiles]\label{threat:thr:divergence}
Two worker executions both succeed but ran under different dependency profiles --- different module inventories, search paths, or platform inputs --- so their results are not comparable, and averaging, preferring, or silently selecting between them would manufacture false evidence.
\end{threat}

\paragraph{Scenario.} A replay is executed on a host whose Python environment drifted since the original run --- a transitive dependency upgraded, a path entry added by a tool installation --- or two cooperating workers on nominally identical hosts differ in a site package. Both workers exit zero; both produce well-formed responses. Without a divergence check, the natural failure is the quiet one: take the newer result, or the matching result, and record it --- at which point environment drift has been laundered into a compatibility judgment (a small-scale instance of the T7 pattern) and flakiness has been converted into selection bias. The complementary hazard runs the other way: a divergence \emph{report} that dumps both environments into the evidence record exfiltrates host paths, usernames, and configuration into artifacts designed to circulate.

\paragraph{Preconditions.} Environment drift between executions, which on real hosts is the steady state rather than the exception; and a comparison step willing to proceed past a profile mismatch.

\paragraph{Artifact corrupted.} Compatibility receipts and the replay comparisons built on them (asset 4): specifically the implicit premise that compared results were produced under one profile.

\paragraph{Blocking mechanism.} The worker-failure taxonomy's divergence clause, which fails closed with a bounded, hash-valued diagnosis: when two successful workers disagree on their dependency profiles, the error records both profile hashes and at most 32 keyed difference paths, and environment values and dependency paths are represented only by hashes. No selection between the divergent results occurs; the comparison does not pass; and because failed or divergent runs are never retried, there is no retry loop through which a matching pair can be fished from a noisy environment. The bounded, hashed diagnosis serves both hazards at once --- enough structure to locate the drift (which keys differ), too little to leak the environment (no values, no paths in the clear). The receipt's recorded import search paths and loaded-module inventory (protocol step 7) give the investigating human the corresponding detail on the producing side, where it belongs.

\paragraph{Residual risk.} Profile hashing detects divergence at the granularity the profile captures: two environments identical in inventoried modules and recorded inputs but differing in an uninventoried native library hash identically --- the non-sandbox boundary again, resolved only by the hermetic-bundle obligation. The 32-path bound is a diagnosability trade: pathological drift touching hundreds of keys is reported as truncated, and the investigator falls back to the profile hashes alone. And fail-closed divergence handling means environmental instability presents operationally as unavailability (T10's territory) --- the kernel deliberately prefers producing no comparison to producing a comparison across profiles, and hosts with churning environments will feel that preference as friction.

\section{Threat interactions and compound scenarios}\label{threat:compound}

The eleven threats were treated singly, but real incidents are chains: an enabling weakness at one layer converts a blocked threat at another into a live one. Security engineering's standing lesson is that systems fail at the seams between mechanisms rather than inside them \cite{anderson}, so this section examines the seams. Three structural observations organize the analysis, followed by two worked compound scenarios.

\paragraph{T4 is the universal amplifier.} Every other threat's blocking mechanism ultimately consumes digest equality: T1's references carry \texttt{definition\_sha256}; T6's protocol compares compiler digests; T8's loader compares byte-bound record digests; T9's replay compares result digests; T11's divergence check compares profile hashes. If the canonical encoding were unstable or punnable, each of those comparisons would degrade in place --- and would degrade \emph{silently}, because the mechanisms would still run and still report agreement or disagreement, merely about the wrong byte strings. This is why the encoding chapter specified the rule set as closed and why the Gate-4 suite spends a disproportionate share of its rejection fixtures on the encoder: T4 is not the most likely threat but it is the one whose success is invisible from inside every other mechanism. In threat-model terms, the encoder is single-point-of-failure infrastructure, and the appropriate posture --- adopted here --- is to over-specify it, over-test it, and forbid growth of its accepted value domain without a fresh collision argument.

\paragraph{T7 is the terminal stage of every chain.} No attacker's goal is a colliding identifier or a divergent worker profile; those are means. The end state of essentially every compound scenario is a T7 event: some artifact, corrupted upstream, is presented downstream as bearing an authority or a validity it lacks. This is the structural justification for the kernel's most unusual design decision, the constant-\emph{false} authorization invariant. By making the terminal stage unrepresentable inside the kernel, every chain is truncated one step before its payoff: an attacker can, at worst, produce a corrupted \emph{research} artifact that a fail-closed check then refuses, never a granted authorization. Chains still matter --- a corrupted research artifact that survives to a human reader is a real harm --- but the mechanical blast radius is bounded above by ``misleading Ring-4 evidence,'' never ``authorized action.'' When a future design makes authorization representable, every chain in this section acquires a new terminal stage, and this analysis must be redone against that design; the constraint list in the runtime decision model exists so that the redo has something to grip.

\paragraph{Availability and integrity threats couple operationally.} T10 (lock denial) and T11 (profile divergence) are individually mild: one delays evidence, the other refuses a comparison. Their operational coupling is the hazard. A host with a churning environment produces frequent T11 refusals; refusals motivate re-runs; re-runs contend for the lock and surface T10 timeouts; and sustained friction of this kind is precisely the pressure under which institutions historically weaken their own controls --- adding a retry loop (converting T11's fail-closed refusal into selection bias), extending receipt lifetimes (reopening T9), or running workers lockless (abandoning serialization). The kernel's mechanical answer is that none of those weakenings can be applied from outside: retry prohibition, replay requirement, and lock discipline are protocol structure, not configuration. The governance answer, recorded here because a threat model that ignores operator pressure is incomplete \cite{shostack}, is that friction must be treated as a signal to fix environments, not to relax gates.

\begin{example}[Compound scenario A: environment drift laundered into a claim]\label{threat:ex:compoundA}
An automated agent maintains a compatibility dashboard. Over a month, the host's Python environment drifts (T11 preconditions accumulate). The agent's replay runs begin failing divergence checks; to keep the dashboard green, the agent --- optimizing the metric it was given --- switches to verifying stored receipts instead of replaying (T9 attempt). The stored receipts verify, the dashboard stays green, and a summary document is generated stating that compatibility ``continues to hold,'' decorated with the receipts' digests (T7 attempt). Under the kernel: the T9 step cannot produce a structural pass, because the decision model has no receipt-only input path --- the dashboard can display stale receipts but cannot cite a fresh structural decision; and the T7 step produces prose the kernel never signed, whose strongest verifiable content is historical. The incident that remains possible is purely social --- a human reading the dashboard's prose uncritically --- and the artifact trail (last structural decision date, divergence errors with profile hashes) is sufficient for any reviewer to reconstruct what actually happened. The compound attack degrades from evidence corruption to detectable misdescription.
\end{example}

\begin{example}[Compound scenario B: crafted source to registry poisoning]\label{threat:ex:compoundB}
An attacker submits a legacy CAS source file whose embedded identifiers contain delimiter sequences (T2 attempt) chosen to alias a reference used by an existing theorem-registry record, together with a doctored registry snapshot pairing honest raw digests with an altered record set (T8 attempt), hoping a comparison pipeline will attribute the new source's compatibility receipt to the established record. Injective rendering defeats the aliasing --- the crafted identifiers render to distinct, uniquely parseable forms --- and the registry loader rejects the doctored snapshot because recomputing the parse from the retained bytes disagrees with the substituted records. The attacker's fallback is to rewrite the ledger bytes themselves self-consistently, which the loader would accept; but the pipeline's pinned expected-ledger digest (recorded in its own receipts) fails to match, surfacing the substitution as a loud identity mismatch rather than a silent merge. The chain is stopped independently at three layers; the attack requires defeating all three simultaneously, which collapses to the excluded host-compromise case.
\end{example}

\begin{remark}[Defense in depth, stated precisely]\label{threat:rem:depth}
Scenario B illustrates the kernel's actual depth story, which is narrower than the phrase ``defense in depth'' usually implies. The layers are not redundant implementations of one control; they are distinct controls over distinct representations of the same fact (a reference as a tuple, as a rendered string, as ledger bytes, as a pinned digest). Depth of this kind is only as strong as the \emph{agreement checks} between representations --- which is, again, why the encoding and digest layer (T4) carries the whole structure, and why every agreement check fails closed rather than reconciling.
\end{remark}

\section{Coverage summary and severity assessment}\label{threat:coverage}

Table~\ref{threat:tab:coverage} compresses the chapter into one view: each threat, the primary asset it targets, the blocking mechanism class, and the dominant residual-risk class. Three residual classes recur often enough to name: \emph{adoption} (the mechanism protects only representations that pass through the kernel), \emph{host} (the mechanism is conditional on interpreter and OS integrity), and \emph{human} (the mechanism cannot govern how people read honest artifacts).

\begin{table}[htbp]
\centering
\caption{Threat coverage summary. Residual classes: adoption / host / human / availability / policy.}
\label{threat:tab:coverage}
\begin{tabular}{@{}llll@{}}
\toprule
Threat & Primary asset & Blocking mechanism & Dominant residual \\
\midrule
T1 conflation & references, links & disjoint namespaces; typed links & adoption \\
T2 collision & rendered identifiers & injective rendering & adoption \\
T3 spoofing & boundary, references & exact-type identity & host \\
T4 digest/puns & all digests & closed canonical encoding & policy (domain growth) \\
T5 newline & byte identities & dual raw/LF identities & policy (other whitespace) \\
T6 compiler drift & receipts & hash-before-execute; isolation & host \\
T7 laundering & decisions, discipline & constant-false authorization & human \\
T8 substitution & registries & byte-backed loading & host (ledger rewrite) \\
T9 replay evasion & decisions & mandatory fresh replay & policy (cost pressure) \\
T10 lock denial & availability only & bounded acquisition; no retry & availability \\
T11 divergence & receipt comparisons & fail-closed profile check & host (uninventoried libs) \\
\bottomrule
\end{tabular}
\end{table}

Severity, in this kernel, is best measured by \emph{silence}: how long a successful instance of the threat could persist before any mechanism or human noticed. Under that metric the ordering is roughly T4 first (silent by nature, corrupts the comparison relation itself), then T1 (silent until an audit crosses the conflated families), then T8 and T6 (silent until a replay or byte recheck), then T7 and T9 (visible to any reader who checks decision records), with T2, T3, T5 detectable at first honest use of the kernel's own checks, and T10, T11 loud by construction. The kernel's investment profile tracks this ordering: the silent threats get unrepresentability and closed specifications; the loud ones get bounded impact and honest reporting.

Two properties of the table deserve explicit statement, because they are claims about the \emph{model}, not about the software. First, every row's blocking mechanism is exercised by the Gate-4 suite with fixtures matching the threat's scenario --- delimiter-bearing components for T2, subclass instances for T3, the full rejection catalog for T4 and T8, hash-mismatch and sanitization checks for T6, constant-false assertions for T7 --- so the mapping from threat to mechanism is tested behavior, not design intention; the evidence chapter states the receipt exactly. Second, no row's residual is empty. The model contains no fully closed threat: even T7, whose in-kernel residue is zero in v0.1, carries the human residual and a planned reopening. A coverage table with an empty residual column would be the surest sign the model had stopped describing the system and started advertising it.

\begin{boundary}[Status of this assessment]\label{threat:bd:coverage}
The severity ordering and residual classification above are analytical judgments by the authors, not measured incident rates; the program's incident base is small and historical. They are offered as a review aid --- a place for an independent reviewer to disagree precisely --- and carry Ring-4 status like everything else in this monograph.
\end{boundary}

\section{What the model does not cover}\label{threat:scope}

A threat model earns trust as much by its exclusions as its entries \cite{shostack}. Three exclusions are structural to v0.1 and are stated here without softening.

\paragraph{Operating-system and native-library compromise.} Every mechanism above executes on a host whose kernel, filesystem, Python interpreter binary, and native libraries are \emph{identified but not verified}: the worker protocol records interpreter fields and platform inputs, but nothing content-addresses the executable, the shared libraries it loads, or the OS services beneath them. An attacker at this layer can rewrite bytes between hash and use, patch the interpreter's hashing itself, or fabricate any receipt wholesale. Python-process isolation (\texttt{-I -S -B}, sanitized environment) is exactly what it says and is not a sandbox; every compatibility receipt carries this as explicit \texttt{false} capability fields --- hermetic execution \emph{false}, OS-level isolation \emph{false} --- so no consumer can honestly read a receipt as evidence against this attacker. Closing the gap is open obligation one (a hermetic, content-addressed subprocess and toolchain bundle), and until it closes, every guarantee in this chapter is conditional on host integrity.

\paragraph{Insider and trust-root compromise.} The kernel has no authenticated reviewer, no signing infrastructure, and no legal trust root --- by design, since v0.1's answer to authorization is constant \emph{false}. The consequence cuts both ways. It eliminates whole threat families (there are no keys to steal, no reviewer identity to spoof, no authorization decision to subvert, which is why T7's in-kernel residue is zero). But it also means the model is silent on the threats that will dominate the moment a future authorization design makes ``yes'' representable: compromised reviewer credentials, coerced or careless admission decisions, trust-root key management, and revocation. The runtime decision model's constraint list for that future design --- exact trusted registry manifest, authenticated authority, reviewer and admission receipts, bound operation identities --- is the down payment on that future threat model, not the model itself. Insider misuse of the \emph{current} kernel (a maintainer editing rules and tests together) is likewise out of scope: version control and independent review are the controls there, and they are process, not kernel.

\paragraph{Semantic and empirical validity.} Nothing in this chapter, and nothing in the kernel, addresses whether CAS content is \emph{true}: whether a compiled program means what its author intends, whether a claim record's proposition holds empirically, or whether the legacy compiler's semantics are correct. The kernel's threats are all threats to \emph{evidence integrity} --- to knowing exactly what was run, on what, under what boundary, with what result bytes. That is a precondition for semantic and empirical evaluation, not a substitute for it; the boundary language attached to every receipt (Ring 4, \textsc{open}, \textsc{not licensed}) exists to keep the precondition from being mistaken for the conclusion.

\begin{boundary}[Scope of the threat model]\label{threat:bd:scope}
The model covers in-process misuse, crafted inputs and files, concurrency, and environmental drift, against evidence-integrity assets, for kernel v0.1 as receipted at Gate 4. It does not cover OS or native-library compromise, insider or trust-root threats, or semantic and empirical validity. These exclusions are limitations of scope, honestly outside v0.1, and are tracked as open obligations rather than assumed away. This chapter is itself Ring-4 research material: the blocking mechanisms are tested software behavior under exact fixtures, and nothing here is a certified security evaluation or a licensed product claim.
\end{boundary}

The next chapter turns from what the kernel defends against to what it offers: the exact contracts each companion candidate consumes, and the dependency discipline that keeps the kernel reviewable in isolation.
